\SourceLeaf{20}
\par
Wir bemerken noch, dass, wenn man für $x,y,z$ ihre durch $p,q$ ausgedrückten Werte einsetzt, die Funktion $W$ notwendig identisch $=0$ sein muss. Daher muss auch das Folgende identisch gelten, also unabhängig von den Werten von $\mathrm{d}p,\mathrm{d}q$.
\[
0=(\lambda T+\mu U+\nu V)\mathrm{d}p+(\lambda'T+\mu'U+\nu'V)\mathrm{d}q
\]
Gleichwertig damit ist
\[
\begin{aligned}
\lambda T+\mu U+\nu V&=0\\
\lambda'T+\mu'U+\nu'V&=0
\end{aligned}
\]
Daraus folgt, dass die Größen $\mu\nu'-\nu\mu'$, $\nu\lambda'-\lambda\nu'$, $\lambda\mu'-\mu\lambda'$ der Reihe nach zu $T,U,V$ beziehungsweise zu den Kosinus der Winkel $QX,QY,QZ$ proportional sind. Das hätte man bereits aus dem Vorangegangenen schließen können, wobei die Vorzeichen noch unbestimmt geblieben wären.

\Article{11}
\par
Von diesen allgemeinen Untersuchungen gehen wir zu den elliptischen Sphäroidkörpern über, für die wir sie unternommen haben. Legt man den Koordinatenursprung in den Mittelpunkt des Körpers und bezeichnet die Halbachsen mit $A,B,C$, so lautet die Gleichung der Oberfläche
\[
\frac{xx}{AA}+\frac{yy}{BB}+\frac{zz}{CC}=1
\]
\par
Wir setzen daher $W=\frac{xx}{AA}+\frac{yy}{BB}+\frac{zz}{CC}-1$. Offenbar nimmt $W$ in allen Punkten innerhalb des Körpers negative und in allen Punkten außerhalb positive Werte an. Ferner gilt $T=\frac{2x}{AA}$, $U=\frac{2y}{BB}$, $V=\frac{2z}{CC}$. Setzen wir also
\[
\BareRoot{(\frac{xx}{A^4}+\frac{yy}{B^4}+\frac{zz}{C^4})}=\psi
\]
so erhalten wir
\[
\cos QX=\frac{x}{\psi AA},\quad\cos QY=\frac{y}{\psi BB},\quad\cos QZ=\frac{z}{\psi CC}
\]
\[
\cos QM=\frac{1}{\psi r}(\frac{(a-x)x}{AA}+\frac{(b-y)y}{BB}+\frac{(c-z)z}{CC})
\]

\Article{12}
\par
Führen wir nun zwei Unbestimmte $p,q$ so ein, dass gilt
\[
\begin{aligned}
x&=A\cos p\\
y&=B\sin p.\cos q\\
z&=C\sin p.\sin q
\end{aligned}
\]
